% Compete then Collaborate — 한국어판 (XeTeX / tectonic 컴파일: tectonic main_ko.tex)
% 폰트: Noto Serif CJK KR. 모델명·기술용어(RLVR/SFT/GRPO/LoRA/MBPP 등)는 영문 유지.
% v0.2 한국어 리라이팅: 번역투 제거, 자연스러운 한국어 학술 문체 적용.
\documentclass[11pt]{article}

\usepackage{fontspec}
\setmainfont{Noto Serif CJK KR}
\setmonofont{DejaVu Sans Mono}[Scale=0.85]
\usepackage[margin=1in]{geometry}
\usepackage{booktabs}
\usepackage{amsmath,amssymb}
\usepackage{graphicx}
\usepackage{xcolor}
\usepackage{pgfplots}
\pgfplotsset{compat=1.18}
\usepackage[numbers,sort&compress]{natbib}
\usepackage[hidelinks]{hyperref}
\usepackage{kotex}  % 한국어 조사·줄바꿈 처리
\linespread{1.25}

\title{경쟁 후 협력: 프런티어 AI 교사가 만드는 검증 가능한 커리큘럼으로\\
코딩 학생 모델을 모방의 한계 너머로 끌어올리다\\[4pt]
\large\normalfont (Compete then Collaborate)}
\author{김미성(Shawn Kim) \\ \small Genesis Cortex AI Inc. \\ \small ORCID: \href{https://orcid.org/0009-0003-7566-3212}{0009-0003-7566-3212}}
\date{2026년 7월}

\begin{document}
\maketitle

\begin{center}
\small\textbf{코드·데이터·테스트·검증 하네스 공개 저장소:}
\url{https://github.com/shawnkim678/compete-then-collaborate}\\[-2pt]
\small\emph{본문의 모든 수치는 공개 산출물을 이용해 재계산하거나, 각 공급자의 이용약관에 따라 교사 풀이를 재생성함으로써 검증할 수 있다.}
\end{center}

\begin{abstract}
대형 언어모델(LLM)은 이제 소형 \emph{학생(student)} 모델의 학습 데이터를 생성하는 \emph{교사(teacher)}로 활발히 활용된다. 그러나 기존의 다중 교사 지식증류(knowledge distillation) 연구는 여러 교사의 출력을 병합하는 데 그칠 뿐, \textbf{어떤} 프런티어 모델이 실제로 더 잘 가르치는지 묻지 않았다. 게다가 이 물음에 답하려는 시도조차 자기 출력을 선호하는 편향이 알려진 LLM 심판에 의존하는 경우가 많다.

본 연구는 이 두 한계를 정면으로 다루는 \textbf{경쟁 후 협력(compete-then-collaborate)} 프레임워크를 제안한다. 프레임워크는 두 단계로 구성된다. 먼저 네 곳의 대표 연구소에서 개발한 프런티어 AI 교사 --- Claude(Anthropic), Codex-GPT(OpenAI), Grok(xAI), Gemini(Google) --- 를 실행 기반 심판(단위 테스트와 stdin--stdout 대조)과 세 가지 공정성 장치(공유 태스크뱅크, 교사 자기교정, 교집합 통제 학습셋) 아래에서 정면 비교한다. 이어 이렇게 검증된 교사들의 풀이를 모아 하나의 코딩 학생 모델(Qwen2.5-Coder)을 위한 \emph{검증 가능한 커리큘럼}을 구축한다.

실험은 세 가지 결과를 낳는다. 첫째, 실행검증 아래에서 네 교사 모두 자기교정 후에는 표준 문제를 거의 완벽하게 해결한다(약 99--100\%). 이는 실력 차이가 아니라 벤치마크 포화(saturation)에 가깝다. 반면 난도가 높아지면 서열이 뚜렷해진다: Gemini 77\% $>$ Claude 69\% $\approx$ Codex 69\% $>$ Grok 50\%. 다만 뒤에서 보이듯 가장 견고한 결과는 학생 쪽에서 나오며 교사 서열에 좌우되지 않는다.

둘째, 교사의 검증된 풀이를 그대로 \textbf{모방(SFT)}시키면 이미 유능한 코딩 학생 모델은 오히려 \textbf{저하}된다. 7B와 32B 학생 모두에서 이 저하가 일관되게 관찰된다(예: MBPP-test 76.7\%에서 72.7\%로, 경쟁 문제는 5.9\%에서 2.9\%로 --- 전체 교사 합집합 기준).

셋째, \emph{동일한} 협력 커리큘럼을 이번에는 모방 대상이 아니라 \textbf{검증가능보상 강화학습(RLVR)}의 환경으로 사용하면 학생은 \textbf{향상}된다. 홀드아웃 경쟁 문제에서 1000 스텝 학습 동안 5.9\%에서 8.8\%로 정점을 찍는다(상대 향상 +49\%). 곧 SFT의 방향이 뒤집힌다.

요컨대 AI 교사 협력의 가치는 답을 모아 모방시키는 데 있지 않다. 학생이 스스로 실행하며 배우는 \emph{검증 가능한 환경}을 함께 구축하는 데 있다. NVIDIA GB10 상에서 GRPO를 안정적으로 구동하기 위한 프레임워크 패치를 포함해, 온프레미스에서 완전히 재현 가능한 파이프라인을 함께 공개한다.
\end{abstract}

\section{서론}
유능한 ``교사'' LLM을 소형 ``학생'' 모델로 증류하는 작업은 이제 표준 관행이 되었다. 그러나 이 관행에서 두 가지 물음은 여전히 충분히 다뤄지지 않고 있다.

첫째, \emph{어떤} 상용 프런티어 모델이 더 나은 교사인가? 이 물음은 LLM 심판이 부여하는 점수가 아니라 학생 모델의 실제 하위 성능으로 답해야 마땅하다. 둘째, 여러 교사를 \emph{결합}하는 최선의 방식은 무엇인가? 답을 병합하는 것 이외의 다른 기제가 존재하지 않는가?

본 연구는 파이썬 코딩 영역에서 이 두 물음을 함께 다룬다. 우리가 채택한 심판은 \textbf{코드 실행}이다. 실행 결과는 객관적이며, LLM 심판이 보이는 자기선호(self-preference) 편향에서 자유롭다. 실험은 두 단계로 흐른다. 먼저 \textbf{경쟁} 단계에서 교사들은 동일한 태스크뱅크를 푼다. 교사의 출력은 히든 테스트를 통과할 때에만 학생 데이터로 들어간다. 이어 \textbf{협력} 단계에서 교사들의 검증된 결과가 하나의 커리큘럼을 이루며, 이 커리큘럼은 두 갈래로 활용된다 --- 모방 목표로 쓰는 SFT, 그리고 검증가능보상의 강화학습 환경으로 쓰는 RLVR.

\paragraph{본 논문의 기여.}
본 연구는 다음 세 가지를 기여한다. 첫째, 세 가지 공정성 장치(공유 태스크, 자기교정, 교집합) 아래에서 실행검증으로 얻은 편향 없는 \textbf{프런티어 AI 교사 서열}을 제시한다. 둘째, 모방 기반 SFT가 유능한 코딩 학생을 오히려 저하시키는 반면 검증가능보상 강화학습은 향상시킨다는 사실을 통제된 조건 아래에서 대비한 \textbf{협력 방식의 정면 비교}를 제시한다. 셋째, transformers 5.5 / torch 2.11 / cu130 환경에서 GRPO를 구동하기 위해 필요한 프레임워크 패치까지 포함한 \textbf{온프레미스(GB10) 재현 가능 파이프라인}을 공개한다.

\section{관련 연구}
\paragraph{다중 교사 지식증류.}
선행 연구는 여러 교사의 근거(rationale)를 학생 모델로 병합하는 방식을 취해 왔다. 그러나 교사를 무작정 늘리면 \emph{지식 충돌(knowledge conflict)}이 발생해 오히려 성능이 떨어진다는 점, 그래서 정제와 통합이 필요하다는 점이 보고되어 있다\citep{multiteacherkd}. 우리 실험의 합집합-SFT 결과 역시 이 저하 현상을 독립적으로 재현한다. 다만 관점은 다르다. 우리는 모방을 위해 병합하지 않는다. 대신 검증 가능한 환경을 함께 구축한다.

\paragraph{교사 선택과 LLM 심판.}
여러 연구가 견고성 확보를 위해 Claude와 GPT를 ``강한 교사''로 혼용해 왔으며, 동시에 GPT-4가 평가자로 쓰일 때 자기 생성물을 선호한다는 사실이 경고되어 왔다\citep{rlhfbook}. 본 연구는 이 편향을 원천적으로 제거한다 --- 심판은 LLM이 아니라 실행이다.

\paragraph{온-폴리시 증류와 RLVR.}
온-폴리시 증류\citep{minillm}와 대조 증류\citep{distillm2}는 모방 성능을 개선한다. 한편 검증가능보상 강화학습은 최근 추론 모델의 기반이 되고 있다. 본 연구는 이 두 흐름을 연결한다. 교사들이 협력적으로 검증한 문제 자체가 곧 RLVR의 보상 환경이 되는 방식이다.

\paragraph{가상·합성 데이터 기반 앙상블 학습의 역사적 계보.}
\emph{가상(virtual)} 데이터 --- 즉 합성 생성된 데이터 --- 로 학습한 앙상블이 개별 구성원을 능가할 수 있다는 직관은 딥러닝 시대 이전까지 거슬러 올라간다. \citet{jang1999ola}은 1999년 POSTECH에서 가상 데이터 기반 앙상블 학습 알고리즘을 제안했다. 본 연구의 구도는 이 오래된 직관의 현대적 구현이다. 앙상블 구성원의 자리에 프런티어 LLM \emph{교사들}이 놓이고, ``가상 데이터'' 자리에는 실행으로 검증된 합성 문제의 커리큘럼이 놓인다.

물론 이는 방법론적 계승이 아니라 정신적 계보의 연결이다. 1999년의 가상 데이터가 단일 예측기의 앙상블 다양성을 높이는 재료였다면, 본 연구의 검증된 문제는 정책 최적화(RLVR)를 위한 \emph{보상 지형(reward landscape)}을 정의한다. 이 지점에서 본 연구는 앞선 계보의 정신적 후예로 자리한다 --- 앙상블 출력을 평균하는 대신 협력적으로 검증된 커리큘럼을 강화학습 환경으로 사용한다는 것이 그 전환점이다.

\paragraph{연구의 공백.}
지금까지 실행검증으로 \emph{명명된 프런티어 모델을 교사로서 서열화}하고, \emph{모방과 검증가능보상 협력}을 정면 대비하며, 이를 \emph{경쟁 후 협력} 파이프라인으로 묶어낸 연구는 확인하지 못했다.

\section{방법}
\subsection{실행검증 기반 생성 --- 심판의 구성}
각 교사는 과제 하나를 받는다. 과제는 함수 시그니처의 형태이거나 경쟁 문제의 형태이다. 교사는 추론과 함께 하나의 코드 블록을 반환한다. 우리는 여기서 코드만 추출해 격리된 서브프로세스에서 실행한다. 이때 자원 제한과 타임아웃을 걸고, 함수 과제는 단위 테스트 assert로, 경쟁 과제는 stdin/stdout 대조로 판정한다. \textbf{히든 테스트}를 통과한 풀이만 채택한다. 이 설계 덕분에 정답성은 교사 모방이 아니라 실행에서 나온다.

\subsection{공정성 장치를 갖춘 경쟁}
경쟁 단계에는 세 가지 공정성 장치를 둔다.
\begin{itemize}
  \item \textbf{공유 태스크뱅크.} 모든 교사가 동일한 문제를 푼다. 함수 시그니처를 명시해 모호성을 없애고, 외부 의존이 있는 과제는 함수 단위 비교의 타당성을 위해 제외한다.
  \item \textbf{자기교정.} 실패한 과제는 실패 테스트의 오류 메시지와 함께 최대 두 번까지 교사에게 재제시한다. 이는 답을 수정할 기회가 있는 실제 교사의 상황을 모사한다. 이로써 각 교사의 커버리지를 끌어올리고 ``한 번의 운''이 만드는 교란을 제거한다.
  \item \textbf{교집합 통제.} 학생 학습셋은 \emph{모든} 교사가 푼 문제로 제한한다. 학생이 받는 문제와 개수를 동일하게 맞추는 것이다. 결과적으로 학생이 받는 것은 오로지 풀이·설명의 \emph{스타일} 차이뿐이다.
\end{itemize}

\subsection{협력 --- 한 학생을 위한 두 가지 방식}
경쟁 단계를 마친 뒤, 모든 교사의 검증된 풀이의 \textbf{합집합}이 협력 커리큘럼을 이룬다. 이 커리큘럼을 우리는 두 방식으로 활용한다.
\begin{itemize}
  \item \textbf{SFT.} 합집합 풀이를 그대로 모방 대상으로 삼아 학생을 지도학습한다.
  \item \textbf{RLVR.} 검증된 문제를 GRPO 보상 환경으로 삼아, 보상을 통과 테스트 비율로 정의해 학생을 강화학습시킨다.
\end{itemize}
학생 모델은 Qwen2.5-Coder이다. 주 실험은 7B에서, 보조 실험은 32B에서 진행하며 모두 LoRA로 훈련한다.

\section{실험 설정}
\begin{itemize}
  \item \textbf{학생.} Qwen2.5-Coder-7B와 -32B. LoRA(bf16)로 훈련. 하드웨어는 GB10 128GB 통합메모리.
  \item \textbf{교사.} Claude, Codex(GPT), Grok, Gemini. 모두 헤드리스 CLI를 통해 호출한다(네 곳의 대표 연구소인 Anthropic, OpenAI, xAI, Google이 개발). Gemini는 실행검증 \emph{경쟁} 서열에 참여하되, \emph{협력} 실험(SFT/RLVR)은 앞의 세 교사가 검증한 풀이의 합집합만 사용한다.
  \item \textbf{태스크뱅크.} MBPP의 함수 과제(교습 split, MBPP-\emph{test}는 홀드아웃), MBPP 참조 풀이의 변이(mutation)로 만든 버그픽스 과제, 그리고 \texttt{deepmind/code\_contests} 경쟁 문제(난이도 6--9).
  \item \textbf{홀드아웃 평가.} MBPP-test 150문항과 서로소인 경쟁 셋 68문항을 사용한다. 사전에 누수 없음을 확인했다. 지표는 실행 pass@1.
  \item \textbf{RLVR.} TRL GRPOTrainer와 PEFT를 HF 경로에서 사용한다(LoRA 자료형 버그로 Unsloth 커널은 회피). 보상은 테스트 통과 비율에 소량의 형식 보너스를 더한 값으로 정의한다. 롤아웃은 HF-generate 방식이다(현재 스택에서는 vLLM이 호환되지 않는다).
\end{itemize}

\section{결과}
\subsection{교사 경쟁 --- 생성 pass@1}
\paragraph{쉬운 문제(MBPP, 공유 200문항) --- 포화 상태.}
\begin{table}[h]\centering
\begin{tabular}{lcc}
\toprule
교사 & one-shot & 자기교정 후 \\
\midrule
Claude & 96.5\% & \textbf{100\%} \\
Grok   & 89.5\% & 99.5\% \\
Codex  & 90.0\% & 99.0\% \\
Gemini & 85.0\% & 99.5\% \\
\bottomrule
\end{tabular}
\caption{MBPP는 이미 포화 상태이다. 네 교사 모두 자기교정 후 99--100\%에 도달하는데, 이는 실력 차이라기보다 벤치마크 노출을 반영할 가능성이 크다.}
\end{table}

\paragraph{어려운 문제(\texttt{code\_contests}, 난이도 6--9, 150문항) --- 정보량 있는 신호.}
\begin{table}[h]\centering
\begin{tabular}{lcc}
\toprule
교사 & pass@1 (정답) & 코드추출 성공률 \\
\midrule
Gemini & \textbf{77\%} (115/150)$^{\ddagger}$ & 100\%$^{\ddagger}$ \\
Claude & 69\% (104/150)\textsuperscript{\S} & 96\% (144/150) \\
Codex  & 69\% (103/150)          & 90\% \\
Grok   & 50\% (75/150)$^{\dagger}$ & 73\% (109/150)$^{\dagger}$ \\
\bottomrule
\end{tabular}
\caption{어려운 경쟁 문제에서 교사들의 실력이 갈린다: Gemini $>$ Claude $\approx$ Codex $>$ Grok. Gemini가 약 8\%p 앞서며, Claude와 Codex는 한 문항 차이 안에 있다.}
\end{table}

\paragraph{해석(v0.2, 동료검토 반영).}
MBPP에서 관찰된 천장 근접 점수(99--100\%)는 실력 차이가 아니라 \emph{벤치마크 포화 내지 학습 노출}을 반영할 가능성이 크다 --- MBPP는 2021년에 공개된 벤치마크로 네 교사 모두가 학습 시점에 접했을 것이다. 따라서 우리는 MBPP에 근거해 교사 서열을 논하지 \textbf{않는다}.

교사들의 실제 차이는 더 어려운 \texttt{code\_contests}에서 드러난다. \textbf{Gemini가 77\%(115/150)로 앞서고, 그 뒤를 Claude(69\%)와 Codex(69\%)가 사실상 같은 자리에서 잇는다. Grok(50\%)이 가장 뒤에 놓인다.} 흥미로운 점은 어느 특정 공급자의 모델이 어려운 문제에서 유일한 최강자가 아니라는 사실이다. Gemini가 약 8\%p 앞서 있으나 Claude와 Codex 사이 격차는 한 문항 안이다. 무엇보다 이 서열은 어떤 LLM 심판이 아니라 실행 결과만으로 산출되므로, \emph{자기선호나 계열(family) 편향은 설계상 배제}된다. 남은 위험 요소는 교사별로 비대칭적인 학습누수와 CLI/파서 아티팩트이며, 이는 \ref{sec:limits}절에서 다룬다.

\paragraph{공정성 보정.}
실험 도중 우리는 실력 차이와 무관한 세 가지 \emph{인프라} 아티팩트를 확인해 보정했다. 어느 것도 특정 교사에만 유리하도록 다루지 않았으며, 각각 투명하게 보고한다.
\begin{itemize}
  \item \textbf{\S\ Claude.} 첫 실행이 122/150에서 중단되었다. 남은 28문항을 마저 실행(22문항 통과)해 모든 교사를 동일한 150문항 기준으로 채점했다. 그 결과 Claude의 점수는 중단 실행 시의 67\%(82/122)에서 공정한 \textbf{69\%(104/150)}로 재산출되었다.
  \item \textbf{\ddag\ Gemini.} 생성 도중 API의 선불 크레딧이 소진되어 마지막 34/150문항이 과금 오류(429)로 무응답 처리되었다. 크레딧 복구 후 그 34문항만 다시 실행해 병합했다. 최종 점수는 \textbf{77\%(115/150)}, 코드추출 100\%이다.
  \item \textbf{\dag\ Grok.} 초기에는 배치 CLI 호출의 52\%가 빈 응답이었다(오답이 아니라 타임아웃 및 불안정 문제). 빈 응답 재시도 로직을 추가해 다시 측정한 결과 \textbf{50\%(75/150)}가 되었으며, 코드추출도 48\%에서 73\%(109/150)로 회복되었다. 다만 재시도 후에도 27\%(41/150)는 여전히 사용 가능한 코드를 내지 못했다. CLI 불안정과 실제 난이도를 완전히 분리할 수 없기 때문에 이 27\%를 폐기하지 않고 그대로 반영해 보고한다. 그 결과 Grok과 Codex·Claude 사이의 격차(50\% 대 69\%)는 Grok에 유리한 쪽으로 보수적으로 추정된 값이다.
\end{itemize}

\subsection{협력 방식 A --- 모방(SFT)은 유능한 학생에게 도움이 되지 않는다}
홀드아웃 pass@1 (교집합 통제, 공유 197문항):
\begin{table}[h]\centering
\begin{tabular}{lcccc}
\toprule
학생 & base & $+$Claude & $+$Codex & $+$Grok \\
\midrule
7B (MBPP)  & 76.7\% & 74.0\% & 70.0\% & 69.3\% \\
32B (MBPP) & 82.0\% & 80.0\% & 77.3\% & 79.3\% \\
\bottomrule
\end{tabular}
\caption{모든 SFT 학생이 base \emph{아래로} 떨어진다. 다만 교사 서열은 그대로 보존된다(Claude가 가장 덜 해롭다).}
\end{table}

모든 학생이 \textbf{base 아래로} 떨어졌고, 교사 서열은 그대로 보존되었다(Claude가 가장 덜 해로웠다). 전체 교사의 \textbf{합집합}으로 SFT를 해도 결과는 여전히 base 이하이다(MBPP는 72.7\%, 경쟁 문제는 base 5.9\%에서 \textbf{2.9\%}로).

이 규모의 모방은 이미 유능한 코딩 모델을 오히려 저하시킨다. 병목은 모델 크기가 아니라 과제의 \emph{난이도와 개선 여지}에 있다 --- 7B와 32B 모두에서 저하가 확인되기 때문이다.

\subsection{협력 방식 B --- 검증가능보상 강화학습(RLVR)은 학생을 향상시킨다}
\label{sec:rlvr}
동일한 커리큘럼을 이번에는 강화학습 환경으로 사용한다. 경쟁 홀드아웃에서 base가 약하므로 개선 여지가 크다.
\begin{table}[h]\centering
\begin{tabular}{lcc}
\toprule
방식 & 경쟁 base $\rightarrow$ 학생 & 방향 \\
\midrule
SFT (합집합)              & 5.9\% $\rightarrow$ 2.9\%          & $\downarrow$ 저하 \\
RLVR (GRPO, 200 스텝)     & 5.9\% $\rightarrow$ 7.4\%          & $\uparrow$ 상대 +25\% \\
\textbf{RLVR (GRPO, v2, 1000 스텝)} & 5.9\% $\rightarrow$ \textbf{8.8\%} 정점 (1000 스텝에서 7.4\%) & $\uparrow$ \textbf{정점 상대 +49\%} \\
\bottomrule
\end{tabular}
\caption{동일한 데이터, 정반대의 방향: 모방은 학생을 저하시키고 검증가능보상 강화학습은 학생을 향상시킨다. v2의 1000 스텝 실행에서 250--750 스텝 구간에 8.8\% 정점이 형성되며, 후반부에서 7.4\%로 완만히 되돌아간다.}
\end{table}

학습 보상은 초기의 약 0에서 후반의 0.25--1.0까지 상승했고(학습셋에서의 일반화), 홀드아웃 성능 역시 함께 개선되었다. \textbf{동일한 데이터를 어떻게 다루느냐에 따라 방향이 반대로 뒤집힌다.}

\begin{figure}[h]\centering
\begin{tikzpicture}
\begin{axis}[
    name=ax, width=0.85\linewidth, height=6cm,
    xlabel={GRPO 스텝}, xmin=0, xmax=1000,
    axis y line*=left, ylabel={학습 보상 (평균)}, ymin=0, ymax=0.55,
    grid=both, tick align=outside,
    legend style={at={(0.5,-0.30)}, anchor=north, legend columns=2},
]
\addplot[thick, blue, mark=*] coordinates
  {(0,0.217)(100,0.308)(200,0.373)(300,0.339)(400,0.394)(500,0.428)(600,0.371)(700,0.436)(800,0.417)(900,0.418)(1000,0.411)};
\addlegendentry{학습 보상 (평균, 좌축)}
\addlegendimage{thick, red, dashed, mark=square*}
\addlegendentry{held-out pass@1 (우축)}
\end{axis}
\begin{axis}[
    width=0.85\linewidth, height=6cm,
    xmin=0, xmax=1000, hide x axis,
    axis y line*=right, ylabel={held-out pass@1 (\%)}, ymin=0, ymax=12,
]
\addplot[thick, red, dashed, mark=square*] coordinates
  {(0,5.9)(100,5.9)(250,8.8)(500,8.8)(750,8.8)(1000,7.4)};
\end{axis}
\end{tikzpicture}
\caption{v2 학습곡선(1000 스텝). 좌축(파랑): 로그된 학습 보상 평균으로 0.22에서 약 0.43까지 상승 후 정체. 우축(빨강): 68개 경쟁 문항에 대한 체크포인트별 held-out pass@1로, base 5.9\%(4/68)에서 시작해 250--750 스텝의 8.8\%(6/68) 정체 구간까지 오른 뒤 1000 스텝에서 7.4\%(5/68)로 완만히 되돌아간다. 정체 구간 내부의 미세한 차이는 68문항 중 $\pm1$문항에 해당하므로 표집 노이즈 범위 안이다. 견고한 신호는 base에서 RLVR로의 이동, 즉 $4\to 5$--$6$/68이다.}
\end{figure}

\section{논의}
결과의 메시지는 명료하다. \textbf{AI 교사 협력의 가치는 답을 모아 모방시키는 데 있지 않다. 학생이 스스로 실행하며 배우는 검증 가능한 환경을 함께 구축하는 데 있다.} 모방이 전달하는 것은 유능한 코더가 이미 가지고 있는 \emph{스타일}일 뿐이다. 반면 RLVR은 검증된 성공에 보상을 부여함으로써 \emph{능력} 자체를 전달한다.

이 관점은 앞서 재현된 다중 교사 KD 저하 현상(지식 충돌)까지 함께 설명한다. 답을 병합하는 방식은 학생의 모방 상한을 넘어설 수 없지만, 검증가능보상에 대한 강화학습은 이 상한을 넘어설 수 있다.

\section{한계}
\label{sec:limits}
\begin{itemize}
  \item \textbf{벤치마크 포화와 학습 누수.} MBPP는 널리 알려진 2021년 벤치마크이므로 교사들의 천장 근접 점수(99--100\%)는 실력 차이보다 학습셋 노출을 반영할 가능성이 크다. 이 때문에 안정적인 서열 논의는 더 어려운 \texttt{code\_contests} 부분집합에 근거한다. 개별 테스트 문항이 특정 교사의 사전학습 코퍼스에 그대로 존재하는지는 검증하지 않았으므로, 만약 교사 간 비대칭 누수가 존재한다면 서열이 편향될 수 있다.

  \item \textbf{프롬프트·파서·CLI의 대칭성.} 모든 교사가 하나의 프롬프트 템플릿과 하나의 코드추출 파서를 공유한다. 실험 도중 Grok의 배치 CLI가 어려운 문제 호출의 52\%에서 빈 응답을 반환하는 아티팩트가 발견되어(오답이 아닌 타임아웃/불안정), 빈 응답 재시도 로직을 추가해 비교를 공정하게 정렬했다. 남은 형식적 우열이 있더라도 그 폭은 공유 파서와 보고된 추출 성공률 범위 안으로 제한된다.

  \item \textbf{교사 서열은 학생 쪽 결과보다 약하다.} 어려운 부분집합에서 Gemini가 약 8\%p 앞서고, Claude와 Codex는 한 문항 차이 안이며, Grok이 그 뒤에 놓인다. Gemini 아래의 세부 순서는 표본 크기를 고려하면 뚜렷하게 분리되지는 않는다. 반면 \textbf{SFT 실패와 RLVR 성공의 방향 반전은 견고하다}. 이 반전은 \emph{학생 모델}의 홀드아웃 성능 변화이므로 교사 서열이나 어떤 LLM 심판에도 의존하지 않는다.

  \item \textbf{RLVR 향상폭의 절대치.} 희소한 경쟁 보상 아래에서 RLVR 향상은 절대치로 보면 완만하다. 1000 스텝 v2 실행은 홀드아웃 pass@1을 base 5.9\%에서 8.8\% 정체(250--750 스텝)까지 끌어올린 뒤 1000 스텝에서 7.4\%로 완만히 되돌아간다. 홀드아웃이 68문항으로 제한되어 있어 정체 구간 내부의 미세한 차이($\pm 1$문항)는 표집 노이즈 안에 있다. 이 때문에 정확한 정점보다 견고한 방향, 즉 base에서 RLVR로의 이동($4 \to 5$--$6$/68)을 보고하고, 마지막 스텝이 아니라 체크포인트 선택 전략을 권장한다. 더 큰 홀드아웃 셋을 확보하면 이 추정을 더 정밀하게 다듬을 수 있을 것이다.

  \item \textbf{일반화 범위.} 학생 계열은 Qwen2.5-Coder 하나이고, 언어는 파이썬 하나이다.

  \item \textbf{윤리와 이용약관.} 본 연구는 \textbf{학술 연구}(벤치마킹 및 방법론 연구)이지, 어떤 공급자와도 경쟁하는 상용 모델을 개발·배포·증류하는 작업이 아니다. 어떤 학습 모델도 제품으로 제공하지 않는다. 이 목적 설정은 본 연구를 Anthropic(경쟁 제품 및 경쟁 모델 학습 금지)과 OpenAI(OpenAI와 경쟁하는 모델 금지)의 \emph{경쟁 범위} 제한 밖에 놓는다. Google Gemini의 약관도 마찬가지로 경쟁 범위에 한정되어 있다. 이에 더해 우리는 (i)~어떤 교사의 원시 출력도 재배포하지 \textbf{않고}, (ii)~교사 출력으로 학습한 모델도 공개하지 \textbf{않으며}, (iii)~핵심 RLVR 결과는 실행 보상과 공개 경쟁 문제만으로, 즉 \textbf{교사 출력 증류 없이} 도출한다. 이는 xAI의 출력 증류 제한도 함께 해소한다. 정답성은 교사 모방이 아니라 결정적 실행으로 심판되며, 공개된 태스크·테스트·하네스를 통해 누구든지 모든 수치를 오프라인으로 재검증하거나(공개 산출물 재검증) 각자의 공급자 약관에 따라 교사 풀이를 재생성함으로써 재현할 수 있다.

  \item \textbf{하드웨어.} 장시간 RLVR 실행 도중 GPU GSP 타임아웃(Xid 119/154)이 발생해 재부팅이 필요했다. 이에 대응해 롤아웃 강도를 완화하고 체크포인트 저장을 촘촘히 하며 Xid 감시를 병행한다.
\end{itemize}

\section{결론}
프런티어 AI는 편향 없는 실행 기반 심판을 통해 교사로서 서열화될 수 있다. 더 중요한 발견은, 이들의 협력이 \emph{강화학습을 위한 검증 가능한 커리큘럼}의 형태로 표현될 때, 모방이 실패하는 지점에서 코딩 학생을 실제로 향상시킨다는 사실이다. 우리는 이를 NVIDIA GB10에서 재현하기 위한 전체 온프레미스 파이프라인과 프레임워크 패치를 함께 공개한다.

\section*{재현성}
모든 코드, 태스크뱅크, 히든 테스트, 결과, 그리고 본 논문은 \url{https://github.com/shawnkim678/compete-then-collaborate}에서 공개한다. 독자는 보고된 모든 수치를 오프라인으로 재검증하거나(\texttt{python reproduce.py --check-banks --selftest}), 자신의 API 키와 각 공급자의 이용약관에 따라 교사 풀이를 재생성할 수 있다. 우리는 어떤 교사의 원시 출력도 재배포하지 않는다(윤리 진술 참조).

\medskip\noindent
관련 스크립트: \texttt{verify\_code.py}, \texttt{verify\_stdio.py}, \texttt{build\_taskbank\_mbpp.py}, \texttt{build\_contests\_bank.py}, \texttt{prof\_run.py}, \texttt{prof\_run\_stdio.py}, \texttt{prof\_retry.py}, \texttt{equalize\_golden.py}, \texttt{eval\_code\_students.py}, \texttt{grpo\_train.py}. 오케스트레이션은 \texttt{grpo\_orchestrator.sh}(setsid). GB10 프레임워크 패치(trl import 플래그, vLLM off, \texttt{PreTrainedModel.warnings\_issued}, Unsloth LoRA 자료형 버그를 회피하기 위한 HF$+$PEFT 경로)는 연구 로그 \S16--17에 문서화되어 있다.

\section*{감사의 글}
가상 데이터 기반 앙상블 학습에 관한 1999년 학위논문\citep{jang1999ola}으로 본 연구의 검증 가능한 커리큘럼 관점에 개념적 계보를 제공해 주신 장민(Min Jang) 박사(포항공과대학교, POSTECH)께 깊이 감사드린다. 장 박사께서는 초안을 검토해 주시며 본 연구의 엔지니어링 기여 --- NVIDIA GB10 플랫폼에서의 스택 수준 디버깅, 네 곳의 상용 LLM CLI 오케스트레이션, 실행 샌드박스 구축 --- 을 시스템의 주된 강점으로 평하시었다. 본 논문에는 어떤 제3자 저작물도 재배포되지 않으며, 해당 학위논문은 서지사항으로만 인용한다.

\bibliographystyle{plainnat}
\bibliography{refs}

\end{document}
